% =====================================================================
% Appendix J --- The pricing door: derivations, guard and accounting
% ---------------------------------------------------------------------
% Job: the formal backing of Sections 6.4-6.5 (subsec:pricing-capacity,
%      subsec:pricing-markup): the effective-technology and markup hooks,
%      the first-order conditions at marginal cost, the uniqueness guard,
%      the Cobb-Douglas price identity, and the external-account and GDP
%      accounting with the rent booked.
% Code: OpenCGE@src/closures/pricing/pricing.jl, OpenCGE@src/core/equilibrium.jl
%      (_mobile_market_demand, _cost_minimizing_labor), OpenCGE@src/core/
%      diagnostics.jl (gdp_components, retained_profit); the gate lives in
%      experiments/run.jl (evaluate_gates). Records: ADR-0029, ADR-0030, ADR-0031.
% =====================================================================
\section{The pricing door: derivations, guard and accounting}\label{app:pricing-door}

\subsection{Two hooks on one cost line}\label{app:pricing-hooks}
Sector $i$ produces with the CES unit cost
\begin{equation}\label{eq:app-unit-cost}
  c_i \;=\; \Bigl[\,(A^{\mathrm{eff}}_i)^{\epsilon-1}\bigl(\alpha_i\,w_i^{1-\epsilon}
     + (1-\alpha_i)\,P_i^{1-\epsilon}\bigr)\Bigr]^{\frac{1}{1-\epsilon}},
\end{equation}
where $P_i$ is the intermediate price index (CES in $\theta$ over the
purchaser prices $p_j$) and $\alpha_i$ the factor share. The two doors act on
this one line. The capacity door replaces $A_i$ by
$A^{\mathrm{eff}}_i = A_i (y_i/\lambda_i)^{-\delta_i \alpha_i}$ and keeps zero profit,
$p_i = c_i$. The factor share in the exponent converts the value-added
elasticity $\delta_i$ into a gross-output one: at fixed input prices a
Hicks-neutral change $d\ln A^{\mathrm{eff}}_i$ raises value added per worker by
$d\ln A^{\mathrm{eff}}_i/\alpha_i$, so value-added productivity responds to own
output with elasticity $-\delta_i$. The markup door keeps $A^{\mathrm{eff}}_i = A_i$ and sets
$p_i = \mu_i c_i$ with $\mu_i = 1 + \kappa(y_i/\lambda_i - 1)$. Both are
written in $y_i/\lambda_i$, so the reference point ($y = \lambda$) is a
fixed point of either hook ($A^{\mathrm{eff}} = A$, $\mu = 1$) and the
scale-determinacy guard of the kernel, evaluated at the reference, is
unchanged. Their product
$p_i = \bigl(1 + \kappa(y_i/\lambda_i - 1)\bigr)\,c_i\bigl(A_i (y_i/\lambda_i)^{-\delta_i\alpha_i}\bigr)$
is one family with two parameters; $\delta = \kappa = 0$ reproduces the
kernel exactly (bit for bit, as the nesting cells
\texttt{pricing\_capacity-v2-GAMMA-*-d0} and
\texttt{pricing\_markup-v1-GAMMA-*-k0} confirm against
\texttt{matrix\_5x3-v10-GAMMA-*}).

\subsection{First-order conditions at marginal cost}\label{app:pricing-foc}
A firm with a markup buys at the purchaser prices but equates marginal revenue
product to factor cost at its \emph{marginal cost} $c_i = p_i/\mu_i$. Cost
minimisation (Shephard's lemma on \eqref{eq:app-unit-cost}) gives labour demand
\begin{equation}\label{eq:app-labour}
  L_i \;=\; c_i^{\epsilon}\,(A^{\mathrm{eff}}_i)^{\epsilon-1}\,\alpha_i\,y_i\,w_i^{-\epsilon},
\end{equation}
and the demand for input $j$ of user $u$ is
\begin{equation}\label{eq:app-intermediate}
  x_{ju} \;=\; p_j^{-\theta}\,\Omega_{uj}\;
     c_u^{\epsilon}\,(A^{\mathrm{eff}}_u)^{\epsilon-1}\,P_u^{\theta-\epsilon}\,a_u\,y_u ,
\end{equation}
with $a_u$ the domestic intermediate bill per unit of output. The outer
$p_j^{-\theta}$ is the price of the demanded input and stays the purchaser
price; the inner $c_u^{\epsilon}$ is the user's marginal cost. The same factor
$k_u = c_u^{\epsilon}(A^{\mathrm{eff}}_u)^{\epsilon-1}P_u^{1-\epsilon}$ values
the two leaks of the purchaser-price intermediate bill, imported
intermediates $M^{\mathrm{int}}$ and product taxes on intermediate use
$T^{\mathrm{int}}$. These are the four sites where the price enters a
first-order condition (labour demand, the inner factor of intermediate demand,
and both leak valuations); writing $p$ instead of $c$ at any of them inflates
the corresponding demand by $\mu^{\epsilon}$ while leaving the price response
untouched. The two leak sites are invisible on any synthetic fixture
($M^{\mathrm{int}} = T^{\mathrm{int}} = 0$), which is why the paper-side tests
assert the external account on the German table. Under the capacity door
$\mu \equiv 1$, so $c = p$ and only $A^{\mathrm{eff}}$ moves: the capacity door
needs no first-order-condition change.

\subsection{The Cobb--Douglas price identity}\label{app:pricing-cd}
At $\theta = \epsilon = 1$ and a fixed wage $w = 1$ the price block is
log-linear: $\log p_i = \log \mu_i - \log A^{\mathrm{eff}}_i
+ (1-\alpha_i)\sum_j \Omega_{ij}\log p_j$, hence
\begin{equation}\label{eq:app-cd}
  \log p \;=\; (I - B)^{-1}\bigl(\log \mu - \log A^{\mathrm{eff}}\bigr),
  \qquad B = \operatorname{diag}(1-\alpha)\,\Omega .
\end{equation}
A uniform constant markup $\bar\mu$ therefore raises prices by
$\bar\mu^{(I-B)^{-1}\mathbf 1}$, not by $\bar\mu$: the markup cascades through
the intermediate network, by a factor of about $1.8$ on the German table. The
kernel tests assert \eqref{eq:app-cd} for both doors.

\subsection{The uniqueness guard}\label{app:pricing-guard}
With $\delta_i < 0$ a sector's unit cost falls as its output rises, a positive
feedback that can admit several equilibria. Every priced cell is reached by a
deterministic continuation in its door parameter (four equal steps from the
door-free solution) and must pass a uniqueness gate: re-solved from four
starting points (the default initial point, the door-free solution, and the
solution with all quantities scaled by $0.9$ and $1.1$), at least two must
converge and every converged start must lie within $10^{-6}$ of the solution,
and the Jacobian of the full residual must have full rank. The measured
spreads are of order $10^{-9}$ at condition numbers of order $10^{7}$; a
distinct equilibrium would sit about $10^{-2}$ away. All 42 priced cells pass.

\subsection{Accounting with the rent booked}\label{app:pricing-accounting}
The rent is $\Pi = \sum_i (1 - 1/\mu_i)\,p_i y_i$. With retention rate
$s_\pi$ the household spends
$E = (1-\tau)\,w\textstyle\sum_i L_i + F + (1-s_\pi)\Pi$, and the
external-account identity of the kernel becomes
\begin{equation}\label{eq:app-account}
  sE + s_\pi \Pi + T^{\mathrm{int}} + M - (I + X) \;=\; F + B_{\mathrm{gov}},
\end{equation}
which every priced run satisfies to $10^{-11}$; dropping the $s_\pi\Pi$ term
leaves it short by exactly the retained rent. On the income side GDP is factor
income $w\sum_i L_i + \Pi$, so the companion identity between the GDP wedge and
the account gap holds unchanged. Retention ($s_\pi = 1$) needs injections in
the fixed-wage gauge: summing over sectors,
$(\mu - 1)\,c\,y + (s + m - s m)\,w L = p\,(1-m)(g^G + \mathit{exo}) + pX$, so an
economy without government demand, investment or exports has no positive
retained-rent equilibrium. The application carries all three.
